\chapter{Pipelines and automation}
\label{ch:pipelines}

\begin{quote}
\emph{``If you do it twice, automate it.''}
--- Every software engineer
\end{quote}

% ============================================================
\section{Why pipelines?}

Manual preprocessing is error-prone and leaky. Common mistakes include:
\begin{itemize}
  \item Fitting a scaler on the full dataset (including test data) --- data leakage.
  \item Forgetting to apply the same transformations at prediction time.
  \item Applying transformations in a different order during training and inference.
\end{itemize}

Scikit-learn's \texttt{Pipeline} and \texttt{ColumnTransformer} solve all three problems by packaging the entire preprocessing workflow into a single, reproducible object.

\begin{datatip}
A pipeline guarantees that every transformation applied during training is applied identically during prediction, in the same order, with the same fitted parameters.
\end{datatip}

% ============================================================
\section{Basic Pipeline}

\begin{minted}{python}
import pandas as pd
import numpy as np
from sklearn.pipeline import Pipeline
from sklearn.preprocessing import StandardScaler
from sklearn.impute import SimpleImputer
from sklearn.linear_model import LogisticRegression
from sklearn.model_selection import train_test_split

# Adult Census dataset
url = ("https://archive.ics.uci.edu/ml/machine-learning-databases/"
       "adult/adult.data")
cols = ["age", "workclass", "fnlwgt", "education", "education_num",
        "marital_status", "occupation", "relationship", "race",
        "sex", "capital_gain", "capital_loss", "hours_per_week",
        "native_country", "income"]
df = pd.read_csv(url, header=None, names=cols, na_values=" ?",
                 skipinitialspace=True)

# Simple numerical pipeline
num_pipe = Pipeline([
    ("imputer", SimpleImputer(strategy="median")),
    ("scaler", StandardScaler()),
])

# Fit and transform
num_cols = ["age", "education_num", "capital_gain",
            "capital_loss", "hours_per_week"]
X_num = num_pipe.fit_transform(df[num_cols])
print(f"Shape: {X_num.shape}")
print(f"Means (should be ~0): {X_num.mean(axis=0).round(2)}")
\end{minted}

% ============================================================
\section{ColumnTransformer}

\begin{minted}{python}
from sklearn.compose import ColumnTransformer
from sklearn.preprocessing import OneHotEncoder

# Define column groups
num_features = ["age", "education_num", "capital_gain",
                "capital_loss", "hours_per_week"]
cat_features = ["workclass", "marital_status", "occupation",
                "relationship", "race", "sex"]

# Build transformers for each type
num_transformer = Pipeline([
    ("imputer", SimpleImputer(strategy="median")),
    ("scaler", StandardScaler()),
])

cat_transformer = Pipeline([
    ("imputer", SimpleImputer(strategy="most_frequent")),
    ("encoder", OneHotEncoder(handle_unknown="ignore",
                               sparse_output=False)),
])

# Combine with ColumnTransformer
preprocessor = ColumnTransformer([
    ("num", num_transformer, num_features),
    ("cat", cat_transformer, cat_features),
])

# Fit and transform
X = preprocessor.fit_transform(df)
print(f"Input shape: {df[num_features + cat_features].shape}")
print(f"Output shape: {X.shape}")
\end{minted}

\begin{preprocesstip}
Use \texttt{ColumnTransformer} to apply different transformations to different column types in a single step. This is the standard approach in production ML systems.
\end{preprocesstip}

% ============================================================
\section{Full pipeline with a model}

\begin{minted}{python}
from sklearn.metrics import accuracy_score, classification_report

# Prepare target
y = (df["income"] == ">50K").astype(int)

# Full pipeline: preprocessing + model
full_pipeline = Pipeline([
    ("preprocessor", preprocessor),
    ("classifier", LogisticRegression(max_iter=1000, random_state=42)),
])

# Train/test split
X_train, X_test, y_train, y_test = train_test_split(
    df[num_features + cat_features], y,
    test_size=0.2, random_state=42, stratify=y
)

# Fit the entire pipeline
full_pipeline.fit(X_train, y_train)

# Predict
y_pred = full_pipeline.predict(X_test)
print(f"Accuracy: {accuracy_score(y_test, y_pred):.4f}")
print(classification_report(y_test, y_pred))
\end{minted}

\begin{minted}{python}
# Cross-validation with the pipeline (no leakage!)
from sklearn.model_selection import cross_val_score

scores = cross_val_score(full_pipeline,
                         df[num_features + cat_features], y,
                         cv=5, scoring="accuracy")
print(f"CV Accuracy: {scores.mean():.4f} +/- {scores.std():.4f}")
\end{minted}

\begin{warningbox}
If you scale your data \emph{before} the train/test split, information from the test set leaks into the scaler's mean and standard deviation. By placing the scaler inside a pipeline, scikit-learn ensures it only fits on training data during cross-validation.
\end{warningbox}

% ============================================================
\section{Custom transformers}

\begin{minted}{python}
from sklearn.base import BaseEstimator, TransformerMixin

class MissingIndicator(BaseEstimator, TransformerMixin):
    """Add binary columns indicating which values were missing."""

    def fit(self, X, y=None):
        self.columns_ = X.columns if hasattr(X, "columns") else None
        return self

    def transform(self, X):
        X = pd.DataFrame(X, columns=self.columns_)
        indicators = X.isnull().astype(int)
        indicators.columns = [f"{c}_missing" for c in indicators.columns]
        return pd.concat([X, indicators], axis=1)


class LogTransformer(BaseEstimator, TransformerMixin):
    """Apply log1p transformation to specified columns."""

    def __init__(self, columns=None):
        self.columns = columns

    def fit(self, X, y=None):
        return self

    def transform(self, X):
        X = X.copy()
        cols = self.columns or X.columns
        for col in cols:
            X[col] = np.log1p(X[col].clip(lower=0))
        return X
\end{minted}

\begin{minted}{python}
# Use custom transformers in a pipeline
custom_pipe = Pipeline([
    ("missing_ind", MissingIndicator()),
    ("imputer", SimpleImputer(strategy="median")),
    ("log", LogTransformer(columns=["capital_gain", "capital_loss"])),
    ("scaler", StandardScaler()),
])

# Test it
sample = df[num_features].head(20)
result = custom_pipe.fit_transform(sample)
print(f"Input columns: {len(num_features)}")
print(f"Output columns: {result.shape[1]}")
\end{minted}

% ============================================================
\section{Saving and loading pipelines}

\begin{minted}{python}
import joblib

# Save the fitted pipeline
joblib.dump(full_pipeline, "income_pipeline.joblib")
print("Pipeline saved.")

# Load it later (e.g., in a web app or batch job)
loaded_pipeline = joblib.load("income_pipeline.joblib")

# Predict on new data
new_data = pd.DataFrame({
    "age": [35], "education_num": [13],
    "capital_gain": [0], "capital_loss": [0],
    "hours_per_week": [40], "workclass": ["Private"],
    "marital_status": ["Married-civ-spouse"],
    "occupation": ["Exec-managerial"],
    "relationship": ["Husband"], "race": ["White"],
    "sex": ["Male"],
})

prediction = loaded_pipeline.predict(new_data)
print(f"Prediction: {'> 50K' if prediction[0] else '<= 50K'}")
\end{minted}

\begin{preprocesstip}
A saved pipeline contains all fitted parameters (scaler means, encoder categories, imputer statistics). This single file is everything you need to preprocess new data identically to training data.
\end{preprocesstip}

% ============================================================
\section{Melbourne Housing full pipeline}

\begin{minted}{python}
# Complete pipeline for the Melbourne Housing dataset
melb = pd.read_csv("melb_data.csv")
melb = melb.dropna(subset=["Price"])

num_features_melb = ["Rooms", "Distance", "Landsize", "BuildingArea",
                     "YearBuilt", "Car", "Bathroom"]
cat_features_melb = ["Type", "Method", "Regionname"]

melb_preprocessor = ColumnTransformer([
    ("num", Pipeline([
        ("imputer", SimpleImputer(strategy="median")),
        ("scaler", StandardScaler()),
    ]), num_features_melb),
    ("cat", Pipeline([
        ("imputer", SimpleImputer(strategy="most_frequent")),
        ("encoder", OneHotEncoder(handle_unknown="ignore",
                                   sparse_output=False)),
    ]), cat_features_melb),
])

from sklearn.ensemble import RandomForestRegressor

melb_pipeline = Pipeline([
    ("preprocessor", melb_preprocessor),
    ("model", RandomForestRegressor(n_estimators=100, random_state=42)),
])

X_melb = melb[num_features_melb + cat_features_melb]
y_melb = melb["Price"]

X_tr, X_te, y_tr, y_te = train_test_split(X_melb, y_melb,
                                           test_size=0.2, random_state=42)
melb_pipeline.fit(X_tr, y_tr)
score = melb_pipeline.score(X_te, y_te)
print(f"Melbourne Housing R^2: {score:.4f}")
\end{minted}

% ============================================================
\section{Exercises}

\begin{exercise}
\begin{enumerate}
  \item Build a \texttt{ColumnTransformer} pipeline for the Titanic dataset with: median imputation + scaling for numerical features, and mode imputation + one-hot encoding for categorical features. Evaluate with 5-fold cross-validation.
  \item Create a custom transformer \texttt{OutlierClipper} that clips values at the 1st and 99th percentiles. Integrate it into a pipeline.
  \item Build a complete pipeline for the Melbourne Housing dataset that includes feature engineering (price per room, building age). Save it with \texttt{joblib} and reload it to make a prediction.
  \item Compare pipeline accuracy with and without preprocessing on the Adult Census dataset. Use a decision tree classifier.
  \item Create a pipeline that handles numerical, categorical, and text columns simultaneously using \texttt{ColumnTransformer}. Use TF-IDF for the text column.
\end{enumerate}
\end{exercise}

% ============================================================
\section{Chapter summary}

\begin{itemize}
  \item Scikit-learn \texttt{Pipeline} chains transformations and a model into a single object, preventing data leakage.
  \item \texttt{ColumnTransformer} applies different transformations to different column types.
  \item Custom transformers (inheriting from \texttt{BaseEstimator} and \texttt{TransformerMixin}) integrate seamlessly into pipelines.
  \item Saving pipelines with \texttt{joblib} captures all fitted parameters for reproducible deployment.
  \item A pipeline ensures that training-time and prediction-time preprocessing are identical.
\end{itemize}
